% Opening the Black Box: Mechanistic Interpretability of Agent Tool Selection
% with Sparse Autoencoders --- Conference poster
% Built on the gemini beamerposter template (https://github.com/anishathalye/gemini)

\documentclass[final]{beamer}

% ====================
% Packages
% ====================

\usepackage[T1]{fontenc}
\usepackage{lmodern}
\usepackage[size=custom,width=140,height=108,scale=1.05]{beamerposter}
\usetheme{gemini}
\usecolortheme{forest}
\usepackage{graphicx}
\usepackage{booktabs}
\usepackage{tikz}
\usepackage{pgfplots}
\pgfplotsset{compat=1.14}
\usepackage{anyfontsize}
\usepackage{qrcode}
\usepackage{xcolor}

% Boost header contrast: white title on dark green
\setbeamercolor{headline}{fg=white}
\setbeamercolor{headline title}{fg=white}
\setbeamercolor{headline author}{fg=white}
\setbeamercolor{headline institute}{fg=white!85}

% ====================
% Brand typography --- Source Sans Pro (Dataiku Designer runtime font)
% Overrides the gemini theme's Lato/Raleway with a single brand family so
% the poster does not mix font systems (DESIGN.md). Must come after
% \usetheme{gemini}, which defines \Lato and \Raleway via \newfontfamily.
% ====================
\setsansfont{Source Sans Pro}[
  Ligatures=TeX,
  UprightFont=*-Regular,
  ItalicFont=*-It,
  BoldFont=*-Semibold,
  BoldItalicFont=*-SemiboldIt
]
\renewfontfamily\Raleway{Source Sans Pro}[Ligatures=TeX,BoldFont=*-Semibold]
\renewfontfamily\Lato{Source Sans Pro}[
  Ligatures=TeX,
  UprightFont=*-Regular,
  ItalicFont=*-It,
  BoldFont=*-Semibold,
  BoldItalicFont=*-SemiboldIt
]

% ====================
% Lengths
% ====================

\newlength{\sepwidth}
\newlength{\colwidth}
\setlength{\sepwidth}{0.025\paperwidth}
\setlength{\colwidth}{0.3\paperwidth}

\newcommand{\separatorcolumn}{\begin{column}{\sepwidth}\end{column}}

% ====================
% Title
% ====================

\title{Opening the Black Box: Mechanistic Interpretability of Agent Tool Selection with Sparse Autoencoders}

\author{Hannes Hapke \inst{1} \and David Cardozo \inst{1}}

\institute[Dataiku 575 Lab]{\inst{1} 575 Lab, Dataiku Inc.\ --- Portland, OR \& Toronto, ON}

% ====================
% Logos
% ====================

\logoright{
\includegraphics[height=5.5cm]{images/575lab_white_condensed.png}}

% ====================
% Footer
% ====================

\footercontent{
  \href{https://github.com/dataiku/kiji-inspector}{github.com/dataiku/kiji-inspector} \hfill
  Mechanistic Interpretability for AI Agents \hfill
  \href{mailto:hannes.hapke@dataiku.com}{hannes.hapke@dataiku.com} \,$\cdot$\,
  \href{mailto:david.cardozo@dataiku.com}{david.cardozo@dataiku.com}}

% ====================
% Body
% ====================

\begin{document}

\begin{frame}[t]
\begin{columns}[t]
\separatorcolumn

% ============================================================
% COLUMN 1 --- Hero result + motivation + pipeline
% ============================================================
\begin{column}{\colwidth}

  \begin{alertblock}{Key Result --- Causal, Not Just Correlated}
    \begin{center}
      \vspace{-0.1em}
      {\fontsize{180}{195}\selectfont\textbf{10.1\,\%}}\\[0.2em]
      {\LARGE of agent tool choices \textbf{reverse}}\\[0.1em]
      {\LARGE when 10 SAE features are ablated.}\\[0.5em]
      {\large \emph{Removing a handful of features
                    \textbf{causally changes} the agent's tool choice.}}\\[0.5em]
      {\large random ablation: \textbf{0\,\%} \quad$\cdot$\quad
              SAE round-trip: \textbf{0\,\%} \quad$\cdot$\quad
              $p = 0.002$}\\[0.3em]
      {\normalsize \emph{up to 17.1\,\% on other contrast types.}}\\[0.4em]
      {\footnotesize \emph{fundamental vs.\ technical analysis} (finance
       domain, 1 of 23 contrast types tested) ---
       Nemotron-3-Nano-30B, layer 20.}
      \vspace{-0.1em}
    \end{center}
  \end{alertblock}

  \begin{block}{The Tool-Selection Black Box}

    \large
    Modern AI agents autonomously pick among databases, code, web search, and
    sub-agents. When the agent chooses, \emph{why} it chose remains opaque.

    \begin{itemize}
      \item \textbf{Prompt engineering:} correlational, not causal.
      \item \textbf{Behavioural testing:} input/output only.
      \item \textbf{Chain-of-thought:} plausible, often confabulated.
    \end{itemize}

    We open the black box at the token where the model commits to a tool and
    decompose its hidden state into human-readable features.

  \end{block}

  \begin{block}{Pipeline}

    \begin{figure}
      \centering
      \begin{tikzpicture}[
        every node/.style={align=center, font=\normalsize},
        main/.style={rectangle, draw=fdark!70, rounded corners=2pt,
                     line width=1.2pt, minimum height=2.6cm,
                     inner xsep=0.4cm, inner ysep=0.3cm},
        arrow/.style={->, line width=1.8pt, >=stealth, draw=fdark},
        sidebox/.style={rectangle, draw=fdark!60, rounded corners=2pt,
                        line width=1.2pt, fill=fbeige,
                        minimum height=3.0cm, minimum width=7cm,
                        inner sep=0.3cm},
        branchArrow/.style={->, line width=1.8pt, >=stealth,
                            draw=fdark!70, dashed},
      ]
        % Main top-row pipeline (purely unsupervised up to "SAE features")
        \node[main, fill=fdark!12,  minimum width=4.0cm] (a) at (0,0)
              {\textbf{Activations}\\\small 1\,M @ layer 20};
        \node[main, fill=fmint!30,  minimum width=5.0cm] (t) at (6.5,0)
              {\textbf{SAE training}\\\small \emph{unsupervised}\\\small 10{,}752 feats};
        \node[main, fill=fmint!30,  minimum width=4.0cm] (s) at (12.5,0)
              {\textbf{SAE features}};
        \node[main, fill=fbeige,    minimum width=4.0cm] (p) at (17.5,0)
              {\textbf{Probe}\\\small Cohen's $d$};
        \node[main, fill=fbeige,    minimum width=4.2cm] (v) at (22.5,0)
              {\textbf{Fuzzing}\\\small token-level A/B};
        \draw[arrow] (a) -- (t);
        \draw[arrow] (t) -- (s);
        \draw[arrow] (s) -- (p);
        \draw[arrow] (p) -- (v);

        % Contrastive pairs branch joins ONLY at the Probe step.
        % "Enters here, not at training" lives INSIDE the side box so it
        % doesn't run into the Fuzzing box.
        \node[sidebox] (c) at (17.5,-5.0)
              {\textbf{Contrastive pairs}\\\small 500K $\cdot$ 5 domains\\[0.2em]
               \small\itshape\textcolor{fdark}{enters here, not at training}};
        \draw[branchArrow] (c) -- (p);
      \end{tikzpicture}
      \caption{\textbf{The SAE is trained without contrastive pairs.}
        Pairs enter only at the probe step to surface which pre-existing
        features carry decision-relevant signal --- never as a training
        signal.}
    \end{figure}

  \end{block}

  \begin{block}{Pick the Right Layer (Not the MoE Ones)}

    \begin{figure}
      \centering
      \includegraphics[width=0.95\linewidth,keepaspectratio]{images/chart_layer_sweep.png}
      \caption{\textbf{A single SAE cannot decompose MoE layers.}
        Reconstruction MSE explodes by $>$500$\times$ at layers 32 and 44
        (sparse expert routing); layer~20 sits just before the MoE stack
        and gives 81.2\,\% alive features. \emph{Don't probe MoE layers
        with one SAE.}}
    \end{figure}

  \end{block}

\end{column}

\separatorcolumn

% ============================================================
% COLUMN 2 --- Methodology + Production hook
% ============================================================
\begin{column}{\colwidth}

  \begin{block}{JumpReLU Sparse Autoencoder}

    \begin{figure}
      \centering
      \includegraphics[width=0.82\linewidth,keepaspectratio]{images/sae_architecture.png}
      \caption{\textbf{An overcomplete dictionary that decomposes hidden states
        into monosemantic features.} JumpReLU thresholds $\theta_i$ produce
        \emph{exact} zeros (true sparsity), trained with a smooth $\tanh$
        surrogate~[2].}
    \end{figure}

    \begin{table}
      \centering
      \large
      \begin{tabular}{ll}
        \toprule
        \textbf{Dictionary $M$}         & 10{,}752 ($4\times$ hidden dim) \\
        \textbf{Activations trained on} & 1{,}000{,}000 \\
        \textbf{Reconstruction MSE}     & 0.574 \\
        \textbf{Alive features}         & 81.2\,\% \\
        \textbf{Active per input}       & $668 \pm 26$ (6.2\,\%) \\
        \bottomrule
      \end{tabular}
    \end{table}

  \end{block}

  \begin{block}{Inside a Feature --- What an SAE Actually Finds}

    \begin{center}
      {\Large\textbf{Feature label:}}\\[0.2em]
      {\LARGE\textbf{\emph{``data modification intent''}}}
    \end{center}

    \vspace{0.4em}
    {\Large\textcolor{fdark}{\textbf{$\blacktriangle$ Activates on}}}\\[0.15em]
    {\large\texttt{set\,$\cdot$\,update\,$\cdot$\,modify\,$\cdot$\,create\,$\cdot$\,delete\,$\cdot$\,write}}

    \vspace{0.5em}
    {\Large\textcolor{fink!55}{\textbf{$\blacktriangledown$ Suppressed by}}}\\[0.15em]
    {\large\texttt{what\,$\cdot$\,list\,$\cdot$\,show\,$\cdot$\,describe\,$\cdot$\,find\,$\cdot$\,latest}}

    \vspace{0.7em}
    \heading{\Large Behaviour under ablation}

    \large
    \begin{tabular}{@{}r@{\,\,}c@{\,\,}l@{}}
      Write request, feature \emph{intact}    & $\Rightarrow$ & agent picks \texttt{[file\_write]}\\[0.15em]
      Write request, feature \emph{ablated}   & $\Rightarrow$ & agent drifts to \texttt{[file\_read]}\\
    \end{tabular}

    \vspace{0.5em}
    {\small Features like this emerge \emph{without} contrastive supervision
    --- the SAE learns them from raw activations alone.}

  \end{block}

  \begin{block}{Scope \& Limitations}
    \begin{itemize}
      \item One model (Nemotron-3-Nano-30B), one layer, one decision-token
            convention.
      \item Top-10 ablation is conservative; distributed decisions need
            larger sets.
      \item Feature labels rely on Qwen3-VL-235B --- judge errors propagate
            to rationales.
    \end{itemize}
  \end{block}

  \begin{exampleblock}{Why This Matters in Production}

    \begin{figure}
      \centering
      \includegraphics[width=0.74\linewidth,keepaspectratio]{images/inference_pipeline.png}
      \caption{\textbf{Every tool call ships with a rationale.}
        At inference, the top-$K$ active features at the decision token
        map to pre-computed labels --- a human-readable explanation
        alongside the tool choice, with no extra LLM call.}
    \end{figure}

    \large
    Most SAE work stops at finding features. Here the features
    \textbf{ship to inference}: every tool call carries a top-$K$ feature
    rationale alongside it --- a \textbf{debugging surface} for deployed
    agents, and a foundation for \textbf{future anomaly detection} when
    activations deviate from a tool's normal fingerprint.

  \end{exampleblock}

\end{column}

\separatorcolumn

% ============================================================
% COLUMN 3 --- Results + Scope + References
% ============================================================
\begin{column}{\colwidth}

  \begin{block}{Feature Label Quality (Token-Level Fuzzing)}

    \begin{figure}
      \centering
      \begin{tikzpicture}
        \colorlet{tierExcellent}{fdark}
        \colorlet{tierGood}{fmint}
        \colorlet{tierPoor}{fbeige}
        % Stacked bar 20cm x 4.2cm
        \fill[tierExcellent] (0,0)       rectangle (16.860,4.2); % 84.3 of 20
        \fill[tierGood]      (16.860,0)  rectangle (18.740,4.2); %  9.4
        \fill[tierPoor]      (18.740,0)  rectangle (20.000,4.2); %  6.2
        \draw[line width=0.6pt] (0,0) rectangle (20,4.2);
        % Hero label inside Excellent (well-separated rows)
        \node[white,font=\fontsize{56}{60}\selectfont\bfseries]
              at (8.43,2.8) {84.3\,\%};
        \node[white,font=\LARGE]
              at (8.43,0.9) {Excellent ($>0.8$)};
        % Leader lines + callouts below for Good / Poor
        \draw[line width=0.6pt] (17.800,0) -- (17.800,-0.6);
        \draw[line width=0.6pt] (19.370,0) -- (19.370,-0.6);
        \node[font=\small,anchor=north east] at (17.800,-0.6) {Good 9.4\,\% ~};
        \node[font=\small,anchor=north west] at (19.370,-0.6) {~ Poor 6.2\,\%};
      \end{tikzpicture}
      \caption{\textbf{84\,\% of feature labels are excellent under
        token-level fuzzing.} 402 features, 4{,}422 A/B examples; combined
        score weights token-level $0.7$ vs.\ prompt-level $0.3$.}
    \end{figure}

    \begin{itemize}
      \item \textbf{Combined score:} $0.912 \pm 0.008$ ($p<10^{-4}$ vs.\ 0.5).
      \item \textbf{Token-level accuracy:} $0.906 \pm 0.007$.
      \item Token-level evaluation catches labels right for the
            \emph{wrong} reasons; prompt-level alone overestimates.
    \end{itemize}

  \end{block}

  \begin{alertblock}{Causal Ablation}

    \begin{figure}
      \centering
      \includegraphics[width=0.92\linewidth,keepaspectratio]{images/chart_ablation.png}
      \caption{\textbf{Ablating SAE features changes tool decisions ---
        random features do not.} Top-10 contrastive features vs.\ random
        features vs.\ SAE round-trip baseline.}
    \end{figure}

    \large
    \heading{Where it works (and where it doesn't)}

    \begin{itemize}
      \item \textbf{Works:} \emph{fundamental} vs.\ \emph{technical} ---
            10.1\,\% flip, controls 0\,\% ($p=0.002$).
      \item \textbf{Works:} \emph{single} vs.\ \emph{multi-tool} ---
            17.1\,\% vs.\ 2.4\,\% ($p=0.029$).
      \item \textbf{Null: 9 of 23 contrast types show zero flip
            rate.}\\ Those decisions are encoded \emph{distributively}
            across many features, not localized in a sparse circuit.
            A known boundary --- and a checklist item for replications.
    \end{itemize}

    \heading{Why the effect is real}
    Random-ablation flip rate \emph{equals} SAE round-trip distortion across
    all 23 contrast types. The contrastive flip is \textbf{above} this noise
    floor --- it must come from the specific features removed, not from
    encode/decode artefacts.

  \end{alertblock}

  \begin{block}{Code \& Paper}

    \begin{minipage}{0.45\linewidth}
      \centering
      \qrcode[height=9cm]{https://github.com/dataiku/kiji-inspector}\\[0.3em]
      {\small\textbf{Scan for paper + code}}
    \end{minipage}\hfill
    \begin{minipage}{0.52\linewidth}
      \footnotesize\texttt{github.com/dataiku/}\\\texttt{kiji-inspector}

      \vspace{0.6em}
      \footnotesize
      Nemotron-3-Nano-30B (subject)\\
      Qwen3-VL-235B (judge)\\
      Compute by NVIDIA

      \vspace{0.4em}
      \scriptsize Prior work: [1] Bricken'23 $\cdot$
      [2] Rajamanoharan'24 (JumpReLU) $\cdot$
      [3] Juang'24 (autointerp).
    \end{minipage}

  \end{block}

\end{column}

\separatorcolumn
\end{columns}
\end{frame}

\end{document}
